% !PRINTED 261
\weberpage{261}
%<<CONT
\weq{10}{\begin{aligned} A &= \left(\frac{u}{v}\right)^r + \left(\frac{v}{u}\right)^r \\ B &= (uv)^s + \left(\frac{2}{n}\right)^{r+s} \frac{2^s}{(uv)^s} , \end{aligned}}
worin $r$, $s$ zwei ganze Zahlen sind, die den Bedingungen
\weq{11}{(n-1)r \equiv 0, \quad (n+1)s \equiv 0 \ (\mathrm{mod}\ 12)}
genügen, die also, wenn $n$ durch 3 teilbar ist, beide durch 3
teilbar sind und
\[ \left(\frac{2}{n}\right) = (-1)^{\frac{n^2-1}{8}} \]
ist.

Es ergeben sich dann folgende zusammengehörige Vertauschungen:
\weq{12}{\begin{aligned}
&\omega, & & A, & & B, \\
&-\frac{1}{\omega}, & & A, & & B, \\
&\omega+2, & & (-1)^{\frac{(n-1)r}{12}} A, & & (-1)^{\frac{(n+1)s}{12}} B, \\
&\frac{\omega-1}{\omega+1}, & & \left(\frac{2}{n}\right)^r A, & & \left(\frac{2}{n}\right)^r B.
\end{aligned}}

Bilden wir nun aus $A$, $B$ eine ganze rationale Funktion mit
numerischen Koeffizienten
\weq{13}{\Phi_{a,c,\partial} = \Sigma\, M_{h,k}\, A^h B^k,}
worin, falls in (12) Vorzeichenänderungen auftreten, die Exponenten
$h$, $k$ so einzurichten sind, daß in allen Gliedern von (13)
die gleichen Vorzeichenveränderungen stattfinden, so wird das
Funktionensystem $\Phi_{a,c,\partial}$ oder wenigstens $\Phi^2_{a,c,\partial}$ der Forderung 1.
des oben aufgestellten Satzes genügen und wir haben, um auch
die Forderung 2. zu befriedigen und so eine Transformationsgleichung zu erhalten, die Koeffizienten $M_{h,k}$ so zu bestimmen,
daß die sämtlichen $\nu$ Werte von $\Phi_{a,c,\partial}$ für $q=0$ endlich bleiben.
Diese Aufgabe vereinfacht sich wesentlich, wenn $n$ keinen quadratischen Teiler hat, und noch mehr, wenn $n$ eine Primzahl ist.

Hat nämlich $n$ keinen quadratischen Teiler, so kann man aus
$\Phi_{a,0,\partial}$ die sämtlichen Werte $\Phi_{a,c,\partial}$ herleiten, durch Vermehrung
von $\omega$ um gewisse ganze Zahlen; wenn also, was unsere Forderung ist, in der Entwickelung von $\Phi_{a,0,\partial}$ % [sic] Source prints a period after 0 here; parallel occurrences have a comma.
nach steigenden Potenzen
von $q$ keine negativen Potenzen vorkommen, so gilt das Gleiche
von sämtlichen $\Phi_{a,c,\partial}$.
